\subsection{Finite Signing Capacity Under Delayed Sources}
\label{sec:eval-capital}

We vary organization CAL backing and authorization, $B\in\{3600,28800\}$,
against a 1- or 3-second source-publication delay. Each member has the
overlapping capacity $\lfloor2B/3\rfloor$, namely 2,400 or 19,200 CAL.
Three wallets execute independent A$\to$B$\to$C units, each with a
100-CAL principal and separate final, user-funded FUEL inputs for both
payments. The child requests certification using its parent's TXCer.
Once the parent certificate is ready, an independent timer publishes
the parent after the selected delay, regardless of the child's progress.
Automatic source resubmission is disabled in every arm to preserve this
delay; there is no top-up, serialized certification, or time-based unlock.
Execution and Bytes grants are sufficient and incur no resource refusals.

Each cell has three fresh runs, with reversed order in the second
repetition. The driver schedules 1,200 units at 20 units/s over 60\,s,
caps unfinished units at 128, records skipped starts when full, and
allows up to 60\,s of drain. The nine services use in-memory committee,
member, gateway and BlockStore state, with NoSync wallets; this experiment
measures signing-capacity behavior rather than durable-storage latency.

\begin{table}[t]\centering\small
\caption{Finite-capital outcomes per 1,200 planned units. Entries are
  medians [minimum, maximum] of three runs. Partial counts distinct
  payment candidates with one or two retained member approvals.}
\label{tab:capital-current}
\setlength{\tabcolsep}{3pt}
\begin{tabular}{@{}rrrr@{}}\toprule
$B$ / delay & Started & Closed & Partial \\\midrule
3,600 / 1 s & 212 [193,227] & 84 [65,99] & 63 [62,68] \\
3,600 / 3 s & 171 [166,208] & 43 [38,83] & 61 [55,64] \\
28,800 / 1 s & 1,200 & 1,200 & 0 \\
28,800 / 3 s & 1,200 & 1,200 & 0 \\
\bottomrule\end{tabular}\end{table}

At $B=28800$, all six runs close all 1,200 units without a CAL refusal
or residual partial approval. Maximum sampled member reservation is
5,000--5,200 CAL with a one-second delay and 9,000--9,100 CAL with
three seconds, below the 19,200-CAL quota (Fig.~\ref{fig:capital-current}).
At $B=3600$, all members end with 2,400 CAL reserved and none available.
The stopped-state audit accounts for all 96 retained signatures in
one- or two-signature candidates; their 100-CAL reservations exactly
explain the four overlapping local totals. These are distinct input
instances, so no successful input-conflict transition reclaims them.
Between 973 and 1,034 planned units never start, and 125--128 started
units remain unfinished at the bounded stop.

Every publicly executed payment nevertheless closes, all replicas
agree, and the final public gap is zero. User-funded fees are charged
once, with no organization FUEL debit. Thus closure restores capacity
for completed approvals while leaving unrelated partial approvals
charged. The high-budget runs sustain this offered load; the low-budget
runs demonstrate why coverage safety alone does not ensure continued
admission. Unfinished requests remain censored rather than being
assigned successful latencies. Supplement~S4 records the per-member
reservations, fees, failures and complete time series.

\begin{figure}[t]\centering
\includegraphics[width=\columnwidth]{capital-current.pdf}
\caption{CAL reservation divided by member quota: one-second maxima
across sampled members. Each budget has three run traces (solid, dashed,
dotted). The 60-second admission window is followed by bounded drain;
traces end when their runs stop.}
\label{fig:capital-current}
\end{figure}

